% =====================================================
% 题型：余弦定理与正弦值选择题 (q_tjw_20260606_18_cosine_rule_sine_angle.tex)
% =====================================================
% =====================================================
% 难度：★★ (中档)
% =====================================================
\newcommand{\qTriangleCosineRuleSineStem}{%
  在 \(\triangle ABC\) 中，\(\angle ABC = \dfrac{\pi}{4}\)，\(AB = \sqrt{2}\)，\(BC = 3\)，则 \(\sin \angle BAC =\) \hfill （\quad）%
}

\newcommand{\qTriangleCosineRuleSineOptions}{%
  \mcoptions[2]
    {\(\dfrac{\sqrt{10}}{10}\)}
    {\(\dfrac{\sqrt{10}}{5}\)}
    {\(\dfrac{3\sqrt{10}}{10}\)}
    {\(\dfrac{\sqrt{5}}{5}\)}%
}

\newcommand{\qTriangleCosineRuleSineAnswer}{C}

\newcommand{\qTriangleCosineRuleSineSolution}{%
  设 \(AC=b\)，由余弦定理得
  \[
    b^2
    = AB^2 + BC^2 - 2 \cdot AB \cdot BC \cos B
    = 2 + 9 - 2 \cdot \sqrt{2} \cdot 3 \cdot \frac{\sqrt{2}}{2}
    = 5.
  \]
  所以 \(AC=\sqrt{5}\)。

  由正弦定理得
  \[
    \frac{\sin A}{BC} = \frac{\sin B}{AC},
    \qquad
    \sin A
    = \frac{BC \sin B}{AC}
    = \frac{3 \cdot \frac{\sqrt{2}}{2}}{\sqrt{5}}
    = \frac{3\sqrt{10}}{10}.
  \]
  故选 C。%
}

\newcommand{\qTriangleCosineRuleSineDiagnosis}{%
  （留白待收集学生作答后补充）%
}

\newcommand{\qTriangleCosineRuleSineTeachingNotes}{%
  快讲候选。对应图片中的“已知两边及夹角，先余弦定理求第三边，再正弦定理求角”的基本流程。重点追问学生是否能准确对应 \(AB,BC,AC\) 与角 \(A,B,C\) 的对边关系。%
}


% =====================================================
% 别名兼容：旧课次宏定义
% =====================================================
\newcommand{\qTJWZeroSixEighteenStem}{\qTriangleCosineRuleSineStem}
\newcommand{\qTJWZeroSixEighteenOptions}{\qTriangleCosineRuleSineOptions}
\newcommand{\qTJWZeroSixEighteenAnswer}{\qTriangleCosineRuleSineAnswer}
\newcommand{\qTJWZeroSixEighteenSolution}{\qTriangleCosineRuleSineSolution}
\newcommand{\qTJWZeroSixEighteenDiagnosis}{\qTriangleCosineRuleSineDiagnosis}
\newcommand{\qTJWZeroSixEighteenTeachingNotes}{\qTriangleCosineRuleSineTeachingNotes}
